% v2026a; generated by R/47_OA_Examples.R (restricted stage); manifesto text from the Manifesto Project corpus
{\footnotesize
\begin{longtable}{p{0.96\linewidth}}
\caption{Worked examples: populism dimensions}\label{tab:oa9a} \\
\toprule
\textbf{Anti-elitism} --- Socialist Party (Ireland), election February 2011; manifesto 53230\_201102, chunk 1. Scores: Sonnet 9, GPT-4.1-mini 9, Gemini 9. \\
\quad \emph{Quote (Sonnet 4.6):} ``end exorbitant salaries \& pensions of TDs'' \\
\quad \emph{Context:} Abolish the Senate \& end exorbitant salaries \& pensions of TDs \& Ministers. Pay them the average industrial wage \\
\quad \emph{Justification (populism block):} The manifesto consistently frames working-class people as virtuous victims of corrupt elites --- bankers, speculators, and out-of-touch politicians --- in a stark moral struggle, scoring very high across all three populism sub-dimensions. \\
\midrule
\textbf{People-centrism} --- New Zealand First Party (New Zealand), election November 2011; manifesto 64621\_201111, chunk 1. Scores: Sonnet 9, GPT-4.1-mini 9, Gemini 9. \\
\quad \emph{Quote (Sonnet 4.6):} ``welfare of all the people of New Zealand'' \\
\quad \emph{Context:} NZ First believes that the welfare of all the people of New Zealand should come first. It's in our name. \\
\quad \emph{Justification (populism block):} The speech is saturated with populist rhetoric: it frames politics as a binary struggle between "the people" (workers, pensioners, ordinary New Zealanders) and corrupt "special interests," "fat cats," and political elites who have betrayed the nation, scoring very high on all three populism sub-dimensions. \\
\midrule
\textbf{Manichaean} --- Reform UK (United Kingdom), election July 2024; manifesto 51701\_202407, chunk 1. Scores: Sonnet 8, GPT-4.1-mini 8, Gemini 8. \\
\quad \emph{Quote (Gemini 3.5 Flash):} ``betrayed the British people'' \\
\quad \emph{Context:} They have destroyed trust in our democracy and betrayed the British people. Our country is worse off, both financially and culturally. \\
\quad \emph{Justification (populism block):} The text strongly frames the political establishment as corrupt and out of touch, contrasting them with the betrayed British people. \\
\midrule
\textbf{Populism (overall)} --- New Zealand First Party (New Zealand), election November 2011; manifesto 64621\_201111, chunk 1. Scores: Sonnet 9, GPT-4.1-mini 9, Gemini 9. \\
\quad \emph{Quote (Sonnet 4.6):} ``We stand for the welfare of all the people'' \\
\quad \emph{Context:} We stand for the welfare of all the people, and not just the few. We believe NZ can be great again. \\
\quad \emph{Justification (populism block):} The speech is saturated with populist rhetoric: it frames politics as a binary struggle between "the people" (workers, pensioners, ordinary New Zealanders) and corrupt "special interests," "fat cats," and political elites who have betrayed the nation, scoring very high on all three populism sub-dimensions. \\
\bottomrule
\multicolumn{1}{p{0.96\linewidth}}{\scriptsize For each dimension the chunk of an English-language manifesto (United Kingdom, Ireland, Australia, New Zealand, United States) on which all three models give a verified score of at least 7 with a quote that passes the quote check and agree most closely (smallest range; ties: higher mean, then chunk id). Evidence shown: the model with the shortest context. Quote, context and justification as returned by the model. Manifesto text: Manifesto Project corpus \citep{lehmann2026manifesto}, identified by the manifesto id.} \\
\end{longtable}
}
